% Shared preamble for the GibbsQ manuscript.
% Loaded via % Shared preamble for the GibbsQ manuscript.
% Loaded via % Shared preamble for the GibbsQ manuscript.
% Loaded via % Shared preamble for the GibbsQ manuscript.
% Loaded via \input{preamble} from main.tex.

% Packages
\usepackage[T1]{fontenc}
\usepackage[utf8]{inputenc}
\usepackage{textcomp}
\usepackage{newtxtext}
\usepackage[a4paper, top=2.5cm, bottom=2.5cm, left=2cm, right=2cm]{geometry}
\usepackage{amsmath,amsthm}
\let\openbox\relax
\usepackage{newtxmath}
\usepackage{mathtools}
\usepackage{bm}
\usepackage{mathrsfs}
\usepackage{graphicx}
\usepackage{booktabs}
\usepackage{array}
\usepackage{tabularx}
\usepackage[protrusion=true,expansion=true]{microtype}
\usepackage{caption}
\usepackage{enumitem}
\usepackage{thmtools}
\usepackage{titlesec}
\usepackage{xcolor}
\usepackage{hyperref}

% Figure search path
\graphicspath{
  {data/policy/figures/}
  {data/ablation/figures/}
  {data/critical/figures/}
  {data/generalize/figures/}
  {data/stats/figures/}
  {data/drift/figures/}
  {data/reinforce_check/figures/}
  {data/reinforce_train/figures/}
  {data/sweep/figures/}
  {data/stress/figures/}
}

% Tables and floats
\setcounter{topnumber}{3}
\setcounter{bottomnumber}{2}
\setcounter{totalnumber}{5}
\setcounter{secnumdepth}{3}
\setcounter{tocdepth}{2}
\renewcommand{\topfraction}{0.92}
\renewcommand{\bottomfraction}{0.8}
\renewcommand{\textfraction}{0.06}
\renewcommand{\floatpagefraction}{0.8}
\setlength{\textfloatsep}{9pt plus 2pt minus 2pt}
\setlength{\floatsep}{8pt plus 2pt minus 1pt}
\setlength{\intextsep}{8pt plus 2pt minus 1pt}
\renewcommand{\arraystretch}{1.15}
\newcolumntype{Y}{>{\raggedright\arraybackslash}X}
\newcolumntype{L}[1]{>{\raggedright\arraybackslash}p{#1}}
\allowdisplaybreaks[2]

% --- State-of-the-Art Section Hierarchy (Top-Tier Standard) ---

% Level 1: Major Section (The Anchor)
% Standard: Left-aligned, Small Caps, NO bold, NO horizontal rule
\titleformat{\section}
  {\normalfont\large\scshape}
  {\thesection.}
  {0.5em}
  {}
\titlespacing*{\section}{0pt}{3.25ex plus 1ex minus .2ex}{1.5ex plus .2ex}

% Level 2: Minor Section (The Pivot)
% Standard: Left-aligned, Bold, Classical separator
\titleformat{\subsection}
  {\normalfont\normalsize\bfseries}
  {\thesubsection.}
  {0.5em}
  {}
\titlespacing*{\subsection}{0pt}{2.25ex plus 1ex minus .2ex}{1ex plus .2ex}

% Level 3: SubHeading (The Detail)
% Standard: Left-aligned, Italicized
\titleformat{\subsubsection}
  {\normalfont\normalsize\itshape}
  {\thesubsubsection.}
  {0.5em}
  {}
\titlespacing*{\subsubsection}{0pt}{1.5ex plus 0.5ex minus .2ex}{1ex plus .2ex}

% Level 4: Micro-Detail (Paragraph)
% Standard: Run-in style, Bold, with auto-appended colon terminator
\titleformat{\paragraph}[runin]
  {\normalfont\normalsize\bfseries}
  {}
  {0pt}
  {}
  [:~]
\titlespacing*{\paragraph}{0pt}{1.5ex plus 0.2ex}{0em}


% Caption and list styling
\captionsetup{
  font=small,
  labelfont=bf,
  labelsep=period,
  skip=6pt
}
\captionsetup[table]{position=top}
\usepackage{etoolbox}
\AtBeginEnvironment{table}{\centering\small\renewcommand{\arraystretch}{1.25}}
\setlist[itemize]{leftmargin=1.5em, itemsep=4pt, topsep=6pt, parsep=3pt, partopsep=0pt}
\setlist[enumerate]{leftmargin=1.7em, itemsep=5pt, topsep=6pt, parsep=3pt, partopsep=0pt}

% Manual replacements for missing packages
\newcommand{\coloneqq}{\mathrel{\mathop:}=}
\newcommand{\eqqcolon}{=\mathrel{\mathop:}}
\newcommand{\cref}[1]{\autoref{#1}}

% Fix autoref names for shared-counter theorem environments
\def\definitionautorefname{Definition}
\def\propositionautorefname{Proposition}
\def\lemmaautorefname{Lemma}
\def\remarkautorefname{Remark}
\def\corollaryautorefname{Corollary}

% Theorem environments
\declaretheoremstyle[
  spaceabove=8pt,
  spacebelow=8pt,
  headfont=\bfseries,
  bodyfont=\itshape,
  headpunct={.},
  postheadspace=0.6em
]{resultstyle}
\declaretheoremstyle[
  spaceabove=8pt,
  spacebelow=8pt,
  headfont=\bfseries,
  bodyfont=\normalfont,
  headpunct={.},
  postheadspace=0.6em
]{definitionstyle}
\declaretheoremstyle[
  spaceabove=6pt,
  spacebelow=6pt,
  headfont=\itshape,
  bodyfont=\normalfont,
  headpunct={.},
  postheadspace=0.55em
]{remarkstyle}
\declaretheorem[style=resultstyle, numberwithin=section, name=Theorem]{theorem}
\declaretheorem[style=resultstyle, sibling=theorem, name=Proposition]{proposition}
\declaretheorem[style=resultstyle, sibling=theorem, name=Lemma]{lemma}
\declaretheorem[style=resultstyle, sibling=theorem, name=Corollary]{corollary}
\declaretheorem[style=definitionstyle, sibling=theorem, name=Definition]{definition}
\declaretheorem[style=remarkstyle, sibling=theorem, name=Remark]{remark}
\renewcommand*{\theHtheorem}{\thesection.\arabic{theorem}}
\renewcommand*{\theHproposition}{\theHtheorem}
\renewcommand*{\theHlemma}{\theHtheorem}
\renewcommand*{\theHcorollary}{\theHtheorem}
\renewcommand*{\theHdefinition}{\theHtheorem}
\renewcommand*{\theHremark}{\theHtheorem}

% Notation macros
% System parameters
\newcommand{\Nservers}{N}
\newcommand{\arrate}{\lambda}
\newcommand{\srvrate}{\mu}
\newcommand{\totcap}{\Lambda}
\newcommand{\loadcond}{\totcap > \arrate}

% Queue state
\newcommand{\Qstate}{Q}
\newcommand{\Qlen}[1]{Q_{#1}}
\newcommand{\Qmin}{Q_{\min}}

% Routing probabilities
\newcommand{\rprob}[1]{p_{#1}}
\newcommand{\softmax}{\mathrm{softmax}}

% Entropy regularization
\newcommand{\temp}{\alpha}
\newcommand{\entropy}{\mathcal{H}}
\newcommand{\KL}[2]{\mathrm{KL}(#1\|#2)}
\newcommand{\logsumexp}{\mathrm{log\text{-}sum\text{-}exp}}

% Lyapunov
\newcommand{\lyap}{V}
\newcommand{\gen}{\mathcal{L}}
\newcommand{\drift}{\gen \lyap}

% UAS-specific
\newcommand{\uaspot}[1]{\Phi_{#1}}
\newcommand{\uasprior}[1]{r_{#1}}
\newcommand{\uasenergy}[1]{x_{#1}}

% Calibrated UAS
\newcommand{\caluas}{\mathrm{cal}}
\newcommand{\scuas}{\mathrm{SC}}
\newcommand{\calbeta}{\beta}
\newcommand{\calgamma}{\gamma}
\newcommand{\caloffset}{c}

% Neural policy
\newcommand{\ngibbsq}{\text{N-GibbsQ}}
\newcommand{\bc}{\text{BC}}
\newcommand{\reinforce}{\text{REINFORCE}}

% Probability and expectation
\newcommand{\E}{\mathbb{E}}
\newcommand{\Prob}{\mathbb{P}}
\newcommand{\Ind}{\mathbf{1}}

% Sets and spaces
\newcommand{\Nats}{\mathbb{Z}_+}
\newcommand{\Reals}{\mathbb{R}}
\newcommand{\Simplex}{\Delta_{\Nservers-1}}

% Norms
\newcommand{\norm}[1]{\left\lVert #1 \right\rVert}
\newcommand{\abs}[1]{\left\lvert #1 \right\rvert}
\newcommand{\onenorm}[1]{\left\lvert #1 \right\rvert_1}

% Operators
\DeclareMathOperator*{\argmin}{arg\,min}
\DeclareMathOperator*{\argmax}{arg\,max}

% Colors for emphasis
\definecolor{resultblue}{RGB}{31, 119, 180}
\definecolor{resultgreen}{RGB}{44, 160, 44}

% Hyperref setup
\hypersetup{
  colorlinks=true,
  linkcolor=blue!70!black,
  citecolor=green!50!black,
  urlcolor=blue!80!black,
  bookmarksnumbered=true,
  pdfstartview={FitH},
}

% Title block
\makeatletter
\renewcommand{\maketitle}{%
  \begin{center}
    {\LARGE\bfseries \@title\par}
    \vspace{0.9em}
    {\large \@author\par}
    \ifx\@date\@empty\else
      \vspace{0.5em}
      {\small \@date\par}
    \fi
  \end{center}
  \vspace{1.2em}
}
\makeatother
 from main.tex.

% Packages
\usepackage[T1]{fontenc}
\usepackage[utf8]{inputenc}
\usepackage{textcomp}
\usepackage{newtxtext}
\usepackage[a4paper, top=2.5cm, bottom=2.5cm, left=2cm, right=2cm]{geometry}
\usepackage{amsmath,amsthm}
\let\openbox\relax
\usepackage{newtxmath}
\usepackage{mathtools}
\usepackage{bm}
\usepackage{mathrsfs}
\usepackage{graphicx}
\usepackage{booktabs}
\usepackage{array}
\usepackage{tabularx}
\usepackage[protrusion=true,expansion=true]{microtype}
\usepackage{caption}
\usepackage{enumitem}
\usepackage{thmtools}
\usepackage{titlesec}
\usepackage{xcolor}
\usepackage{hyperref}

% Figure search path
\graphicspath{
  {data/policy/figures/}
  {data/ablation/figures/}
  {data/critical/figures/}
  {data/generalize/figures/}
  {data/stats/figures/}
  {data/drift/figures/}
  {data/reinforce_check/figures/}
  {data/reinforce_train/figures/}
  {data/sweep/figures/}
  {data/stress/figures/}
}

% Tables and floats
\setcounter{topnumber}{3}
\setcounter{bottomnumber}{2}
\setcounter{totalnumber}{5}
\setcounter{secnumdepth}{3}
\setcounter{tocdepth}{2}
\renewcommand{\topfraction}{0.92}
\renewcommand{\bottomfraction}{0.8}
\renewcommand{\textfraction}{0.06}
\renewcommand{\floatpagefraction}{0.8}
\setlength{\textfloatsep}{9pt plus 2pt minus 2pt}
\setlength{\floatsep}{8pt plus 2pt minus 1pt}
\setlength{\intextsep}{8pt plus 2pt minus 1pt}
\renewcommand{\arraystretch}{1.15}
\newcolumntype{Y}{>{\raggedright\arraybackslash}X}
\newcolumntype{L}[1]{>{\raggedright\arraybackslash}p{#1}}
\allowdisplaybreaks[2]

% --- State-of-the-Art Section Hierarchy (Top-Tier Standard) ---

% Level 1: Major Section (The Anchor)
% Standard: Left-aligned, Small Caps, NO bold, NO horizontal rule
\titleformat{\section}
  {\normalfont\large\scshape}
  {\thesection.}
  {0.5em}
  {}
\titlespacing*{\section}{0pt}{3.25ex plus 1ex minus .2ex}{1.5ex plus .2ex}

% Level 2: Minor Section (The Pivot)
% Standard: Left-aligned, Bold, Classical separator
\titleformat{\subsection}
  {\normalfont\normalsize\bfseries}
  {\thesubsection.}
  {0.5em}
  {}
\titlespacing*{\subsection}{0pt}{2.25ex plus 1ex minus .2ex}{1ex plus .2ex}

% Level 3: SubHeading (The Detail)
% Standard: Left-aligned, Italicized
\titleformat{\subsubsection}
  {\normalfont\normalsize\itshape}
  {\thesubsubsection.}
  {0.5em}
  {}
\titlespacing*{\subsubsection}{0pt}{1.5ex plus 0.5ex minus .2ex}{1ex plus .2ex}

% Level 4: Micro-Detail (Paragraph)
% Standard: Run-in style, Bold, with auto-appended colon terminator
\titleformat{\paragraph}[runin]
  {\normalfont\normalsize\bfseries}
  {}
  {0pt}
  {}
  [:~]
\titlespacing*{\paragraph}{0pt}{1.5ex plus 0.2ex}{0em}


% Caption and list styling
\captionsetup{
  font=small,
  labelfont=bf,
  labelsep=period,
  skip=6pt
}
\captionsetup[table]{position=top}
\usepackage{etoolbox}
\AtBeginEnvironment{table}{\centering\small\renewcommand{\arraystretch}{1.25}}
\setlist[itemize]{leftmargin=1.5em, itemsep=4pt, topsep=6pt, parsep=3pt, partopsep=0pt}
\setlist[enumerate]{leftmargin=1.7em, itemsep=5pt, topsep=6pt, parsep=3pt, partopsep=0pt}

% Manual replacements for missing packages
\newcommand{\coloneqq}{\mathrel{\mathop:}=}
\newcommand{\eqqcolon}{=\mathrel{\mathop:}}
\newcommand{\cref}[1]{\autoref{#1}}

% Fix autoref names for shared-counter theorem environments
\def\definitionautorefname{Definition}
\def\propositionautorefname{Proposition}
\def\lemmaautorefname{Lemma}
\def\remarkautorefname{Remark}
\def\corollaryautorefname{Corollary}

% Theorem environments
\declaretheoremstyle[
  spaceabove=8pt,
  spacebelow=8pt,
  headfont=\bfseries,
  bodyfont=\itshape,
  headpunct={.},
  postheadspace=0.6em
]{resultstyle}
\declaretheoremstyle[
  spaceabove=8pt,
  spacebelow=8pt,
  headfont=\bfseries,
  bodyfont=\normalfont,
  headpunct={.},
  postheadspace=0.6em
]{definitionstyle}
\declaretheoremstyle[
  spaceabove=6pt,
  spacebelow=6pt,
  headfont=\itshape,
  bodyfont=\normalfont,
  headpunct={.},
  postheadspace=0.55em
]{remarkstyle}
\declaretheorem[style=resultstyle, numberwithin=section, name=Theorem]{theorem}
\declaretheorem[style=resultstyle, sibling=theorem, name=Proposition]{proposition}
\declaretheorem[style=resultstyle, sibling=theorem, name=Lemma]{lemma}
\declaretheorem[style=resultstyle, sibling=theorem, name=Corollary]{corollary}
\declaretheorem[style=definitionstyle, sibling=theorem, name=Definition]{definition}
\declaretheorem[style=remarkstyle, sibling=theorem, name=Remark]{remark}
\renewcommand*{\theHtheorem}{\thesection.\arabic{theorem}}
\renewcommand*{\theHproposition}{\theHtheorem}
\renewcommand*{\theHlemma}{\theHtheorem}
\renewcommand*{\theHcorollary}{\theHtheorem}
\renewcommand*{\theHdefinition}{\theHtheorem}
\renewcommand*{\theHremark}{\theHtheorem}

% Notation macros
% System parameters
\newcommand{\Nservers}{N}
\newcommand{\arrate}{\lambda}
\newcommand{\srvrate}{\mu}
\newcommand{\totcap}{\Lambda}
\newcommand{\loadcond}{\totcap > \arrate}

% Queue state
\newcommand{\Qstate}{Q}
\newcommand{\Qlen}[1]{Q_{#1}}
\newcommand{\Qmin}{Q_{\min}}

% Routing probabilities
\newcommand{\rprob}[1]{p_{#1}}
\newcommand{\softmax}{\mathrm{softmax}}

% Entropy regularization
\newcommand{\temp}{\alpha}
\newcommand{\entropy}{\mathcal{H}}
\newcommand{\KL}[2]{\mathrm{KL}(#1\|#2)}
\newcommand{\logsumexp}{\mathrm{log\text{-}sum\text{-}exp}}

% Lyapunov
\newcommand{\lyap}{V}
\newcommand{\gen}{\mathcal{L}}
\newcommand{\drift}{\gen \lyap}

% UAS-specific
\newcommand{\uaspot}[1]{\Phi_{#1}}
\newcommand{\uasprior}[1]{r_{#1}}
\newcommand{\uasenergy}[1]{x_{#1}}

% Calibrated UAS
\newcommand{\caluas}{\mathrm{cal}}
\newcommand{\scuas}{\mathrm{SC}}
\newcommand{\calbeta}{\beta}
\newcommand{\calgamma}{\gamma}
\newcommand{\caloffset}{c}

% Neural policy
\newcommand{\ngibbsq}{\text{N-GibbsQ}}
\newcommand{\bc}{\text{BC}}
\newcommand{\reinforce}{\text{REINFORCE}}

% Probability and expectation
\newcommand{\E}{\mathbb{E}}
\newcommand{\Prob}{\mathbb{P}}
\newcommand{\Ind}{\mathbf{1}}

% Sets and spaces
\newcommand{\Nats}{\mathbb{Z}_+}
\newcommand{\Reals}{\mathbb{R}}
\newcommand{\Simplex}{\Delta_{\Nservers-1}}

% Norms
\newcommand{\norm}[1]{\left\lVert #1 \right\rVert}
\newcommand{\abs}[1]{\left\lvert #1 \right\rvert}
\newcommand{\onenorm}[1]{\left\lvert #1 \right\rvert_1}

% Operators
\DeclareMathOperator*{\argmin}{arg\,min}
\DeclareMathOperator*{\argmax}{arg\,max}

% Colors for emphasis
\definecolor{resultblue}{RGB}{31, 119, 180}
\definecolor{resultgreen}{RGB}{44, 160, 44}

% Hyperref setup
\hypersetup{
  colorlinks=true,
  linkcolor=blue!70!black,
  citecolor=green!50!black,
  urlcolor=blue!80!black,
  bookmarksnumbered=true,
  pdfstartview={FitH},
}

% Title block
\makeatletter
\renewcommand{\maketitle}{%
  \begin{center}
    {\LARGE\bfseries \@title\par}
    \vspace{0.9em}
    {\large \@author\par}
    \ifx\@date\@empty\else
      \vspace{0.5em}
      {\small \@date\par}
    \fi
  \end{center}
  \vspace{1.2em}
}
\makeatother
 from main.tex.

% Packages
\usepackage[T1]{fontenc}
\usepackage[utf8]{inputenc}
\usepackage{textcomp}
\usepackage{newtxtext}
\usepackage[a4paper, top=2.5cm, bottom=2.5cm, left=2cm, right=2cm]{geometry}
\usepackage{amsmath,amsthm}
\let\openbox\relax
\usepackage{newtxmath}
\usepackage{mathtools}
\usepackage{bm}
\usepackage{mathrsfs}
\usepackage{graphicx}
\usepackage{booktabs}
\usepackage{array}
\usepackage{tabularx}
\usepackage[protrusion=true,expansion=true]{microtype}
\usepackage{caption}
\usepackage{enumitem}
\usepackage{thmtools}
\usepackage{titlesec}
\usepackage{xcolor}
\usepackage{hyperref}

% Figure search path
\graphicspath{
  {data/policy/figures/}
  {data/ablation/figures/}
  {data/critical/figures/}
  {data/generalize/figures/}
  {data/stats/figures/}
  {data/drift/figures/}
  {data/reinforce_check/figures/}
  {data/reinforce_train/figures/}
  {data/sweep/figures/}
  {data/stress/figures/}
}

% Tables and floats
\setcounter{topnumber}{3}
\setcounter{bottomnumber}{2}
\setcounter{totalnumber}{5}
\setcounter{secnumdepth}{3}
\setcounter{tocdepth}{2}
\renewcommand{\topfraction}{0.92}
\renewcommand{\bottomfraction}{0.8}
\renewcommand{\textfraction}{0.06}
\renewcommand{\floatpagefraction}{0.8}
\setlength{\textfloatsep}{9pt plus 2pt minus 2pt}
\setlength{\floatsep}{8pt plus 2pt minus 1pt}
\setlength{\intextsep}{8pt plus 2pt minus 1pt}
\renewcommand{\arraystretch}{1.15}
\newcolumntype{Y}{>{\raggedright\arraybackslash}X}
\newcolumntype{L}[1]{>{\raggedright\arraybackslash}p{#1}}
\allowdisplaybreaks[2]

% --- State-of-the-Art Section Hierarchy (Top-Tier Standard) ---

% Level 1: Major Section (The Anchor)
% Standard: Left-aligned, Small Caps, NO bold, NO horizontal rule
\titleformat{\section}
  {\normalfont\large\scshape}
  {\thesection.}
  {0.5em}
  {}
\titlespacing*{\section}{0pt}{3.25ex plus 1ex minus .2ex}{1.5ex plus .2ex}

% Level 2: Minor Section (The Pivot)
% Standard: Left-aligned, Bold, Classical separator
\titleformat{\subsection}
  {\normalfont\normalsize\bfseries}
  {\thesubsection.}
  {0.5em}
  {}
\titlespacing*{\subsection}{0pt}{2.25ex plus 1ex minus .2ex}{1ex plus .2ex}

% Level 3: SubHeading (The Detail)
% Standard: Left-aligned, Italicized
\titleformat{\subsubsection}
  {\normalfont\normalsize\itshape}
  {\thesubsubsection.}
  {0.5em}
  {}
\titlespacing*{\subsubsection}{0pt}{1.5ex plus 0.5ex minus .2ex}{1ex plus .2ex}

% Level 4: Micro-Detail (Paragraph)
% Standard: Run-in style, Bold, with auto-appended colon terminator
\titleformat{\paragraph}[runin]
  {\normalfont\normalsize\bfseries}
  {}
  {0pt}
  {}
  [:~]
\titlespacing*{\paragraph}{0pt}{1.5ex plus 0.2ex}{0em}


% Caption and list styling
\captionsetup{
  font=small,
  labelfont=bf,
  labelsep=period,
  skip=6pt
}
\captionsetup[table]{position=top}
\usepackage{etoolbox}
\AtBeginEnvironment{table}{\centering\small\renewcommand{\arraystretch}{1.25}}
\setlist[itemize]{leftmargin=1.5em, itemsep=4pt, topsep=6pt, parsep=3pt, partopsep=0pt}
\setlist[enumerate]{leftmargin=1.7em, itemsep=5pt, topsep=6pt, parsep=3pt, partopsep=0pt}

% Manual replacements for missing packages
\newcommand{\coloneqq}{\mathrel{\mathop:}=}
\newcommand{\eqqcolon}{=\mathrel{\mathop:}}
\newcommand{\cref}[1]{\autoref{#1}}

% Fix autoref names for shared-counter theorem environments
\def\definitionautorefname{Definition}
\def\propositionautorefname{Proposition}
\def\lemmaautorefname{Lemma}
\def\remarkautorefname{Remark}
\def\corollaryautorefname{Corollary}

% Theorem environments
\declaretheoremstyle[
  spaceabove=8pt,
  spacebelow=8pt,
  headfont=\bfseries,
  bodyfont=\itshape,
  headpunct={.},
  postheadspace=0.6em
]{resultstyle}
\declaretheoremstyle[
  spaceabove=8pt,
  spacebelow=8pt,
  headfont=\bfseries,
  bodyfont=\normalfont,
  headpunct={.},
  postheadspace=0.6em
]{definitionstyle}
\declaretheoremstyle[
  spaceabove=6pt,
  spacebelow=6pt,
  headfont=\itshape,
  bodyfont=\normalfont,
  headpunct={.},
  postheadspace=0.55em
]{remarkstyle}
\declaretheorem[style=resultstyle, numberwithin=section, name=Theorem]{theorem}
\declaretheorem[style=resultstyle, sibling=theorem, name=Proposition]{proposition}
\declaretheorem[style=resultstyle, sibling=theorem, name=Lemma]{lemma}
\declaretheorem[style=resultstyle, sibling=theorem, name=Corollary]{corollary}
\declaretheorem[style=definitionstyle, sibling=theorem, name=Definition]{definition}
\declaretheorem[style=remarkstyle, sibling=theorem, name=Remark]{remark}
\renewcommand*{\theHtheorem}{\thesection.\arabic{theorem}}
\renewcommand*{\theHproposition}{\theHtheorem}
\renewcommand*{\theHlemma}{\theHtheorem}
\renewcommand*{\theHcorollary}{\theHtheorem}
\renewcommand*{\theHdefinition}{\theHtheorem}
\renewcommand*{\theHremark}{\theHtheorem}

% Notation macros
% System parameters
\newcommand{\Nservers}{N}
\newcommand{\arrate}{\lambda}
\newcommand{\srvrate}{\mu}
\newcommand{\totcap}{\Lambda}
\newcommand{\loadcond}{\totcap > \arrate}

% Queue state
\newcommand{\Qstate}{Q}
\newcommand{\Qlen}[1]{Q_{#1}}
\newcommand{\Qmin}{Q_{\min}}

% Routing probabilities
\newcommand{\rprob}[1]{p_{#1}}
\newcommand{\softmax}{\mathrm{softmax}}

% Entropy regularization
\newcommand{\temp}{\alpha}
\newcommand{\entropy}{\mathcal{H}}
\newcommand{\KL}[2]{\mathrm{KL}(#1\|#2)}
\newcommand{\logsumexp}{\mathrm{log\text{-}sum\text{-}exp}}

% Lyapunov
\newcommand{\lyap}{V}
\newcommand{\gen}{\mathcal{L}}
\newcommand{\drift}{\gen \lyap}

% UAS-specific
\newcommand{\uaspot}[1]{\Phi_{#1}}
\newcommand{\uasprior}[1]{r_{#1}}
\newcommand{\uasenergy}[1]{x_{#1}}

% Calibrated UAS
\newcommand{\caluas}{\mathrm{cal}}
\newcommand{\scuas}{\mathrm{SC}}
\newcommand{\calbeta}{\beta}
\newcommand{\calgamma}{\gamma}
\newcommand{\caloffset}{c}

% Neural policy
\newcommand{\ngibbsq}{\text{N-GibbsQ}}
\newcommand{\bc}{\text{BC}}
\newcommand{\reinforce}{\text{REINFORCE}}

% Probability and expectation
\newcommand{\E}{\mathbb{E}}
\newcommand{\Prob}{\mathbb{P}}
\newcommand{\Ind}{\mathbf{1}}

% Sets and spaces
\newcommand{\Nats}{\mathbb{Z}_+}
\newcommand{\Reals}{\mathbb{R}}
\newcommand{\Simplex}{\Delta_{\Nservers-1}}

% Norms
\newcommand{\norm}[1]{\left\lVert #1 \right\rVert}
\newcommand{\abs}[1]{\left\lvert #1 \right\rvert}
\newcommand{\onenorm}[1]{\left\lvert #1 \right\rvert_1}

% Operators
\DeclareMathOperator*{\argmin}{arg\,min}
\DeclareMathOperator*{\argmax}{arg\,max}

% Colors for emphasis
\definecolor{resultblue}{RGB}{31, 119, 180}
\definecolor{resultgreen}{RGB}{44, 160, 44}

% Hyperref setup
\hypersetup{
  colorlinks=true,
  linkcolor=blue!70!black,
  citecolor=green!50!black,
  urlcolor=blue!80!black,
  bookmarksnumbered=true,
  pdfstartview={FitH},
}

% Title block
\makeatletter
\renewcommand{\maketitle}{%
  \begin{center}
    {\LARGE\bfseries \@title\par}
    \vspace{0.9em}
    {\large \@author\par}
    \ifx\@date\@empty\else
      \vspace{0.5em}
      {\small \@date\par}
    \fi
  \end{center}
  \vspace{1.2em}
}
\makeatother
 from main.tex.

% Packages
\usepackage[T1]{fontenc}
\usepackage[utf8]{inputenc}
\usepackage{textcomp}
\usepackage{newtxtext}
\usepackage[a4paper, top=2.5cm, bottom=2.5cm, left=2cm, right=2cm]{geometry}
\usepackage{amsmath,amsthm}
\let\openbox\relax
\usepackage{newtxmath}
\usepackage{mathtools}
\usepackage{bm}
\usepackage{mathrsfs}
\usepackage{graphicx}
\usepackage{booktabs}
\usepackage{array}
\usepackage{tabularx}
\usepackage[protrusion=true,expansion=true]{microtype}
\usepackage{caption}
\usepackage{enumitem}
\usepackage{thmtools}
\usepackage{titlesec}
\usepackage{xcolor}
\usepackage{hyperref}

% Figure search path
\graphicspath{
  {data/policy/figures/}
  {data/ablation/figures/}
  {data/critical/figures/}
  {data/generalize/figures/}
  {data/stats/figures/}
  {data/drift/figures/}
  {data/reinforce_check/figures/}
  {data/reinforce_train/figures/}
  {data/sweep/figures/}
  {data/stress/figures/}
}

% Tables and floats
\setcounter{topnumber}{3}
\setcounter{bottomnumber}{2}
\setcounter{totalnumber}{5}
\setcounter{secnumdepth}{3}
\setcounter{tocdepth}{2}
\renewcommand{\topfraction}{0.92}
\renewcommand{\bottomfraction}{0.8}
\renewcommand{\textfraction}{0.06}
\renewcommand{\floatpagefraction}{0.8}
\setlength{\textfloatsep}{9pt plus 2pt minus 2pt}
\setlength{\floatsep}{8pt plus 2pt minus 1pt}
\setlength{\intextsep}{8pt plus 2pt minus 1pt}
\renewcommand{\arraystretch}{1.15}
\newcolumntype{Y}{>{\raggedright\arraybackslash}X}
\newcolumntype{L}[1]{>{\raggedright\arraybackslash}p{#1}}
\allowdisplaybreaks[2]

% --- State-of-the-Art Section Hierarchy (Top-Tier Standard) ---

% Level 1: Major Section (The Anchor)
% Standard: Left-aligned, Small Caps, NO bold, NO horizontal rule
\titleformat{\section}
  {\normalfont\large\scshape}
  {\thesection.}
  {0.5em}
  {}
\titlespacing*{\section}{0pt}{3.25ex plus 1ex minus .2ex}{1.5ex plus .2ex}

% Level 2: Minor Section (The Pivot)
% Standard: Left-aligned, Bold, Classical separator
\titleformat{\subsection}
  {\normalfont\normalsize\bfseries}
  {\thesubsection.}
  {0.5em}
  {}
\titlespacing*{\subsection}{0pt}{2.25ex plus 1ex minus .2ex}{1ex plus .2ex}

% Level 3: SubHeading (The Detail)
% Standard: Left-aligned, Italicized
\titleformat{\subsubsection}
  {\normalfont\normalsize\itshape}
  {\thesubsubsection.}
  {0.5em}
  {}
\titlespacing*{\subsubsection}{0pt}{1.5ex plus 0.5ex minus .2ex}{1ex plus .2ex}

% Level 4: Micro-Detail (Paragraph)
% Standard: Run-in style, Bold, with auto-appended colon terminator
\titleformat{\paragraph}[runin]
  {\normalfont\normalsize\bfseries}
  {}
  {0pt}
  {}
  [:~]
\titlespacing*{\paragraph}{0pt}{1.5ex plus 0.2ex}{0em}


% Caption and list styling
\captionsetup{
  font=small,
  labelfont=bf,
  labelsep=period,
  skip=6pt
}
\captionsetup[table]{position=top}
\usepackage{etoolbox}
\AtBeginEnvironment{table}{\centering\small\renewcommand{\arraystretch}{1.25}}
\setlist[itemize]{leftmargin=1.5em, itemsep=4pt, topsep=6pt, parsep=3pt, partopsep=0pt}
\setlist[enumerate]{leftmargin=1.7em, itemsep=5pt, topsep=6pt, parsep=3pt, partopsep=0pt}

% Manual replacements for missing packages
\newcommand{\coloneqq}{\mathrel{\mathop:}=}
\newcommand{\eqqcolon}{=\mathrel{\mathop:}}
\newcommand{\cref}[1]{\autoref{#1}}

% Fix autoref names for shared-counter theorem environments
\def\definitionautorefname{Definition}
\def\propositionautorefname{Proposition}
\def\lemmaautorefname{Lemma}
\def\remarkautorefname{Remark}
\def\corollaryautorefname{Corollary}

% Theorem environments
\declaretheoremstyle[
  spaceabove=8pt,
  spacebelow=8pt,
  headfont=\bfseries,
  bodyfont=\itshape,
  headpunct={.},
  postheadspace=0.6em
]{resultstyle}
\declaretheoremstyle[
  spaceabove=8pt,
  spacebelow=8pt,
  headfont=\bfseries,
  bodyfont=\normalfont,
  headpunct={.},
  postheadspace=0.6em
]{definitionstyle}
\declaretheoremstyle[
  spaceabove=6pt,
  spacebelow=6pt,
  headfont=\itshape,
  bodyfont=\normalfont,
  headpunct={.},
  postheadspace=0.55em
]{remarkstyle}
\declaretheorem[style=resultstyle, numberwithin=section, name=Theorem]{theorem}
\declaretheorem[style=resultstyle, sibling=theorem, name=Proposition]{proposition}
\declaretheorem[style=resultstyle, sibling=theorem, name=Lemma]{lemma}
\declaretheorem[style=resultstyle, sibling=theorem, name=Corollary]{corollary}
\declaretheorem[style=definitionstyle, sibling=theorem, name=Definition]{definition}
\declaretheorem[style=remarkstyle, sibling=theorem, name=Remark]{remark}
\renewcommand*{\theHtheorem}{\thesection.\arabic{theorem}}
\renewcommand*{\theHproposition}{\theHtheorem}
\renewcommand*{\theHlemma}{\theHtheorem}
\renewcommand*{\theHcorollary}{\theHtheorem}
\renewcommand*{\theHdefinition}{\theHtheorem}
\renewcommand*{\theHremark}{\theHtheorem}

% Notation macros
% System parameters
\newcommand{\Nservers}{N}
\newcommand{\arrate}{\lambda}
\newcommand{\srvrate}{\mu}
\newcommand{\totcap}{\Lambda}
\newcommand{\loadcond}{\totcap > \arrate}

% Queue state
\newcommand{\Qstate}{Q}
\newcommand{\Qlen}[1]{Q_{#1}}
\newcommand{\Qmin}{Q_{\min}}

% Routing probabilities
\newcommand{\rprob}[1]{p_{#1}}
\newcommand{\softmax}{\mathrm{softmax}}

% Entropy regularization
\newcommand{\temp}{\alpha}
\newcommand{\entropy}{\mathcal{H}}
\newcommand{\KL}[2]{\mathrm{KL}(#1\|#2)}
\newcommand{\logsumexp}{\mathrm{log\text{-}sum\text{-}exp}}

% Lyapunov
\newcommand{\lyap}{V}
\newcommand{\gen}{\mathcal{L}}
\newcommand{\drift}{\gen \lyap}

% UAS-specific
\newcommand{\uaspot}[1]{\Phi_{#1}}
\newcommand{\uasprior}[1]{r_{#1}}
\newcommand{\uasenergy}[1]{x_{#1}}

% Calibrated UAS
\newcommand{\caluas}{\mathrm{cal}}
\newcommand{\scuas}{\mathrm{SC}}
\newcommand{\calbeta}{\beta}
\newcommand{\calgamma}{\gamma}
\newcommand{\caloffset}{c}

% Neural policy
\newcommand{\ngibbsq}{\text{N-GibbsQ}}
\newcommand{\bc}{\text{BC}}
\newcommand{\reinforce}{\text{REINFORCE}}

% Probability and expectation
\newcommand{\E}{\mathbb{E}}
\newcommand{\Prob}{\mathbb{P}}
\newcommand{\Ind}{\mathbf{1}}

% Sets and spaces
\newcommand{\Nats}{\mathbb{Z}_+}
\newcommand{\Reals}{\mathbb{R}}
\newcommand{\Simplex}{\Delta_{\Nservers-1}}

% Norms
\newcommand{\norm}[1]{\left\lVert #1 \right\rVert}
\newcommand{\abs}[1]{\left\lvert #1 \right\rvert}
\newcommand{\onenorm}[1]{\left\lvert #1 \right\rvert_1}

% Operators
\DeclareMathOperator*{\argmin}{arg\,min}
\DeclareMathOperator*{\argmax}{arg\,max}

% Colors for emphasis
\definecolor{resultblue}{RGB}{31, 119, 180}
\definecolor{resultgreen}{RGB}{44, 160, 44}

% Hyperref setup
\hypersetup{
  colorlinks=true,
  linkcolor=blue!70!black,
  citecolor=green!50!black,
  urlcolor=blue!80!black,
  bookmarksnumbered=true,
  pdfstartview={FitH},
}

% Title block
\makeatletter
\renewcommand{\maketitle}{%
  \begin{center}
    {\LARGE\bfseries \@title\par}
    \vspace{0.9em}
    {\large \@author\par}
    \ifx\@date\@empty\else
      \vspace{0.5em}
      {\small \@date\par}
    \fi
  \end{center}
  \vspace{1.2em}
}
\makeatother
